\section{Implementation}
\label{sec:impl}

We optimize the same high-level regions in AVX2 and AVX-512: SHAKE x1/x4,
party-batched field arithmetic, binary matrix--vector multiplication, and the
first-phase AIM3 MPC simulation.  The two backends differ only in register
layout and instruction selection.  Phase~2/3 arithmetic, delta accumulation,
last-share adjustment, and seed-tree traversal remain scalar control flow;
therefore, this is a matched targeted implementation rather than a claim of
full-scheme vectorization.

\subsection{Four-Way SHAKE}
\label{sec:shake}

Commitment and tape generation for parties $p,p+1,p+2,p+3$ are independent.
For Keccak word $j$, we store the corresponding words of four states in one
YMM register:
\begin{equation}
Y_j=\bigl(A_j^{(p)},A_j^{(p+1)},A_j^{(p+2)},A_j^{(p+3)}\bigr),
\qquad 0\le j<25.
\end{equation}
One vector operation then performs the same Keccak step for all four parties.
Each lane absorbs the same domain prefix, salt, and repetition index but a
different party index and seed.  Outputs are squeezed to four separate
commitments and tapes in their original party order.

The AVX2 backend uses 256-bit YMM operations such as VPXOR, VPANDN, shifts,
shuffles, and permutations.  Rotations and the Keccak $\chi$ operation require
multiple instructions, and only 16 architectural YMM registers are available.
The AVX-512 backend deliberately keeps the same 4-way YMM layout: four
64-bit words require only 256 bits.  Its advantage comes from AVX-512VL rather
than a ZMM-wide x8 construction.  The 25 state words reside in
\code{ymm0}--\code{ymm24}; VPROLQ performs each 64-bit rotation, VPTERNLOGQ
combines XOR/Boolean expressions, and mask registers with VMOVDQU8 handle
partial input and output.  The 32-register file reduces state spills.  The x1
path uses the same AVX-512VL instruction set with one active state.

\subsection{Party-Batched Binary-Field Arithmetic}
\label{sec:field}

Field elements retain the official little-endian 64-bit-word representation.
The irreducible polynomials are
$x^{128}+x^7+x^2+x+1$,
$x^{192}+x^7+x^2+x+1$, and
$x^{256}+x^{10}+x^5+x^2+1$.
PCLMULQDQ produces polynomial partial products, which are combined with XOR
and reduced with constants \code{0x87} for 128/192 bits and \code{0x425} for
256 bits.

In AVX2, carry-less multiplication remains an XMM operation.  A
$\gftwo{128}$ element occupies one XMM register, while 192- and 256-bit
elements use two XMM pieces.  We interleave independent PCLMUL streams to hide
latency.  YMM registers are used for matrix accumulation: one YMM holds two
128-bit parties, one padded 192-bit party, or one 256-bit party.

AVX-512 packs four independent 128-bit pieces into the four lanes of a ZMM
register and evaluates them with VPCLMULQDQ.  For $\gftwo{128}$ this directly
packs four parties.  For $\gftwo{192}$, each party is split into a lower
128-bit part and a zero-padded upper 64-bit part; four lower and four upper
parts form two ZMM registers.  The $\gftwo{256}$ representation similarly
uses lower and upper 128-bit ZMM vectors without padding.  The common
multiplier in a multiply-accumulate operation is broadcast to all four lanes.

\begin{table}[t]
\centering
\caption{Register packing in the matched field implementations.}
\label{tab:packing}
\small
\setlength{\tabcolsep}{5pt}
\begin{tabular}{lcc}
\toprule
Field & AVX2 matrix block & AVX-512 matrix block\\
\midrule
$\gftwo{128}$ & 2 parties/YMM; 8 per loop & 4 parties/ZMM; 16 per loop\\
$\gftwo{192}$ & 1 party/YMM; 4 per loop & 1 party/YMM; 8 per loop\\
$\gftwo{256}$ & 1 party/YMM; 4 per loop & 2 parties/ZMM; 8 per loop\\
\bottomrule
\end{tabular}
\end{table}

The scalar AVX2 and AVX-512 field sources are intentionally identical and use
XMM PCLMULQDQ.  The additional AVX-512 width is used in the party-batch
kernels.  This distinction prevents scalar microbenchmarks from being
misinterpreted as a test of ZMM batching.

\subsection{Matrix Accumulation and MPC Scheduling}
\label{sec:mpc}

For an input bit $b$ and a binary matrix row $r$, the linear layer performs
\begin{equation}
  k\leftarrow k\oplus (\operatorname{mask}(b)\wedge r).
  \label{eq:matrix}
\end{equation}
AVX2 creates the mask with a comparison and evaluates Eq.~\eqref{eq:matrix}
with VPAND followed by VPXOR.  AVX-512 uses VPTERNLOGQ with immediate
\code{0x78}, which evaluates the AND-and-XOR expression in one instruction.
The 192-bit backend uses AVX-512VL masked loads and stores so that only three
64-bit words are accessed.  The 256-bit backend uses VPERMUTEXVARQ to
broadcast the currently tested input word inside each two-party ZMM layout.

The original scalar order completes the MPC calculation one party at a time.
We transpose this schedule: all parties first execute the same AIM3 affine
layer, then the next affine layer, and finally the same sequence of Frobenius
squarings.  The batch functions therefore act as a party-wise map of the
official equations.  Only the last party receives the public affine constant
and ciphertext corrections, exactly as in the reference code.

\subsection{Equivalence Boundary and Library Integration}
\label{sec:equiv}

The SIMD transformation preserves SHAKE rate, padding, domain separators,
absorb/finalize/squeeze order, field representation, AIM3 matrices, random-byte
consumption, proof order, and serialization.  SIMD lanes contain independent
parties and never exchange data.  Thus, vectorization changes the evaluation
schedule but not the mathematical function.

All six parameter sets are registered through an OQS-compatible API.  Runtime
selection prefers AVX-512, then AVX2, and finally reference code; an environment
variable can force each backend for testing.  The project vendors its Keccak
backends and builds without an AIMer v2 source or build dependency.  This is a
self-contained experimental integration rather than a patch to the upstream
liboqs repository~\cite{openquantumsafe}.
